\section {Summary of changes to analysis inputs compared with T2K-TN-036 (v4) }
\label{sec:analysis_modifications}

The following changes have been implemented compared with the $\nu_{\mu}$-disappearance analysis 
described in \cite{OurRun1DisappAnalysis}:

\begin{itemize}
\item
  The SuperK flux has been upgraded from 10d(v2) to 10d(v3) \cite{UpdatedFlux}.
  When fitting the reconstructed energy spectrum of the combined Run1+2 single $\mu$-like ring event sample,
  we used a POT-weighted Run1+Run2 flux in order to take into account the different conditions of Run1 and Run2.
\item
  The flux normalisation ($N_{sk}/N_{nd}$) error has been updated. 
  The updated error calculation code (v3.1) was obtained from \cite{UpdatedFluxNormError}.
  Typical errors used in Ref.\ \cite{OurRun1DisappAnalysis} were
  $\sim$9.7\% without oscillations and
  $\sim$6.3\% for oscillations with 
  $sin^{2}2\theta_{23}$ = 1.0 and $|\Delta m^{2}_{23}|$ = 2.32$\times10^{-3}$ $eV^{2}/c^{4}$.
  The updated flux normalisation errors used in this analysis for these two characteristic cases are 
  $\sim$7.3\% and $\sim$4.7\% respectively.
\item
  The flux shape error has been updated, and
  the updated flux fractional covariance matrix (v3.1) was obtained from \cite{UpdatedFluxNormError}.
\item
  An updated ND280 event rate normalisation correction was used, and it was calculated using the 10d(v3) flux.
  The updated ND280 correction factor and energy spectra for various CC-like event categories 
  were obtained from \cite{UpdatedND280Histograms}.
  The ND280 normalisation correction was based on an analysis of the Run1 data only and,
  as before, it is energy independent.
\item
  The ND280 efficiency error was increased from $\pm$5\% to +5.6\% -5.2\%. The statistical error of the
  observed number of CC-like events was also propagated to SuperK.
\item
  The SuperK ring-counting (RC) efficiency errors have been revised \cite{UpdatedSuperkEfficiencyErrors}.
  The CCQE RC error has been revised from 7.6\% (no energy dependence) to
  1.7\% for $E_{reco}<$ 0.4 GeV, 3.5\% for  0.4 GeV $\leq E_{reco} \leq$ 1.1GeV
  and 9.3\% for  $E_{reco}>$ 1.1 GeV. The bin-to-bin correlations in the RC efficiency error in
  the three reconstructed neutrino energy bins have also been taken into account, and
  the covariance matrix is given in Tab.\ \ref{tab:input_uncertainties}.
  For CCnonQE the RC efficiency error has been revised from 25.4\% to 19.9\% (no energy dependence) 
  and for NC the RC efficiency error has been revised from 57.0\% to 48.2\% (no energy dependence).
  All other SuperK efficiency errors were unchanged.
\item
  The SuperK energy scale error has been revised from 1\% to 2.3\% \cite{UpdatedSuperkEnergyScaleError}.
\end{itemize}


A summary of all the sources of systematic error considered in this
analysis is given in Tab.\ \ref{tab:input_uncertainties}; this is an
update of Tab.\ 6 in Ref.\ \cite{OurRun1DisappAnalysis}.
The treatment of these uncertainties is largely unchanged and is
described in detail in Sec.\ 4.1 and Eq.\ 31 of Ref.\ \cite{OurRun1DisappAnalysis}. 
The only technical complication is that the SuperK
CCQE efficiency error (which was 7.8\% independent of energy in
\cite{OurRun1DisappAnalysis}) 
has been split into two parts: these are the RC efficiency error and the
`other' efficiency error (i.e.\ fiducial volume, reduction, OD cut, PID, momentum cut and decay
electron cut). The CCQE RC efficiency
error now has a dependence on the reconstructed
neutrino energy and the bin-to-bin correlations are taken into account.
The oscillation analysis code has been updated to include these changes.
The bin-to-bin correlations in the RC efficiency error are propagated using the Cholesky 
method \cite{NumRecipCholesky}; this method was employed in the earlier analysis to propagate
the correlated systematic errors in the flux shape (see Sec.\ 4.5 of Ref.\ \cite{OurRun1DisappAnalysis}).

%
% Table with systematic parameters
%

\begin{table*}[htb]
\center
\renewcommand{\arraystretch}{1.75} % enlarge line spacing
\begin{tabular}{p{0.12\textwidth}p{0.88\textwidth}}
%
% SuperK
%
\hline
\multicolumn {2}{l} { {\bf SuperK detector uncertainties} } \\
\hline
\parbox{0.12\textwidth}
{
   $f_{\epsilon;r}^{SK}$  
}&
\parbox{0.88\textwidth}
{
    \vspace{0.2cm}
    Combined uncertainty in efficiency due to uncertainties in fiducial volume, reduction chain, 
    OD cut, RC cut, PID cut, momentum cut and decay electron cut \cite{UpdatedSuperkEfficiencyErrors}.
    Applied to corresponding SuperK MC templates:
    \begin{itemize}
     \item CCQE: 
         \begin{itemize}
            \item RC: 1.7\% ($E_{reco}<$ 0.4 GeV), 3.5\% (0.4 GeV $\leq E_{reco} \leq$ 1.1 GeV), 9.3\% ($E_{reco}>$ 1.1 GeV). 
                  Error covariance matrix is
                  $
                   \left[
                     {\begin{array}{ccc}
                       (1.7\%)^2  &   0.59 * 1.7\% * 3.5\%  &  0.33 * 1.7\% * 9.3\% \\
                                  &   (3.5\%)^2             &  0.61 * 3.5\% * 9.3\% \\
                                  &                         &  (9.3\%)^2            \\
                      \end{array} }
                   \right]
                  $
            \item other: 1.7\% (energy-independent)
         \end{itemize}
     \item CCnonQE: 19.9\% RC + 1.7\% other $\sim$20\% (energy-independent)
     \item NC: 48.2\% RC + 100\% PID + 1.7\% other  $\sim$111\% (energy-independent)
     \item $\nu_{e}$ CC: $\sim$100\% (energy-independent)
    \end{itemize}
}\\
\parbox{0.12\textwidth}
{
   $f_{E}^{SK}$  
}&
\parbox{0.88\textwidth}
{
    \vspace{0.2cm}
    Reconstructed energy scale error is 2.3\% \cite{UpdatedSuperkEnergyScaleError}. 
    Applied to SuperK MC templates by scaling bin edges and,
    assuming uniform distribution of events within the 50 MeV bins, calculating the number of events 
    gained from / lost to neighbouring bins.\\
}\\
\hline
%
% nd280 detector systematics
%
\multicolumn {2}{l} { {\bf ND280 detector uncertainties} } \\
\hline
\parbox{0.12\textwidth}
{ 
   $f_{\epsilon}^{ND280}$  
}&
\parbox{0.88\textwidth}
{
    \vspace{0.2cm}
    Efficiency error is +5.6\% -5.2\% \cite{UpdatedND280Histograms}.
    Applied to ND280 MC templates.\\
}\\
\hline
%
% Beam systematics
%
\multicolumn {2}{l} { {\bf Flux uncertainties} } \\
\hline
\parbox{0.12\textwidth}
{ 
   $f_{norm}$ 
}&
\parbox{0.88\textwidth}
{
    \vspace{0.2cm}
    Flux normalization uncertainty.
    This uncertainty, expressed relative to the $N_{SK}/N_{ND280}$ ratio predicted by the MC,
    was calculated for oscillation and other systematic parameters 
    using code provided by the Beam WG \cite{UpdatedFluxNormError}.
    The value obtained for nominal MC and oscillations with parameters
    $sin^{2}2\theta_{23}$ = 1.00 and $|\Delta m^{2}_{23}|$ = 2.32$\times10^{-3}$ $eV^{2}/c^{4}$
    was $\sim$4.7\% ($\sim$7.3\% without oscillations).\\
}\\
\hline
\parbox{0.12\textwidth}
{ 
    ${\bf \rho^{\Phi_{SK}}}$ 
}&
\parbox{0.88\textwidth}
{
    \vspace{0.2cm}
    Flux shape uncertainty. Expressed in terms of true energy bin-to-bin 
    correlation matrix for SuperK $\nu_{\mu}$ flux provided by the Beam WG \cite{UpdatedFluxShapeError}.
    Uncertainty was propagated using the Choleskly method \cite{NumRecipCholesky}.\\
}\\
\hline
%
% Neutrino Generator
%
\multicolumn {2}{l} { {\bf Neutrino Generator (GENIE/NEUT) model uncertainties} } \\
\hline
\parbox{0.12\textwidth}
{ 
    \vspace{0.2cm}
    $f_{\sigma_{CCQE};t}$ 
}&
\parbox{0.88\textwidth}
{
    \vspace{0.2cm}
    CCQE cross section uncertainty. This respresents only the ``residual'' CCQE cross section uncertainty 
    that is not cancelled in the near-to-far extrapolation; this residual uncertainty is due to the 
    different flux energy ranges sampled by the near and far detectors and the different compositions 
    of their nuclear targets. Therefore the uncertainty is applied to the SuperK MC templates only.
    This uncertainty has dependence on true neutrino energy \cite{NIWGErr}.\\
   
}\\
\parbox{0.12\textwidth}
{ 
    \vspace{0.2cm}
    $f_{\sigma_{CC1\pi}/\sigma_{CCQE};t}$ 
}&
\parbox{0.88\textwidth}
{
    \vspace{0.2cm}
    Uncertainty in CC1$\pi$/CCQE cross section ratio:
    30\% for true energy $<$  2 GeV and 20\% for true energy $>$ 2 GeV \cite{NIWGErr}.
    Applied to both corresponding SuperK and ND280 MC templates.\\
}\\
\parbox{0.12\textwidth}
{ 
    \vspace{0.2cm}
    $f_{\sigma_{CCoth}/\sigma_{CCQE};t}$ 
}&
\parbox{0.88\textwidth}
{
    \vspace{0.2cm}
    Uncertainty in CCother/CCQE cross section ratio:
    30\% for true energy $<$ 2 GeV and 25\% for true energy $>$ 2 GeV \cite{NIWGErr}.
    Applied to both corresponding SuperK and ND280 MC templates.\\ 
}\\
\parbox{0.12\textwidth}
{ 
    \vspace{0.2cm}
    $f_{\sigma_{NC}/\sigma_{CCQE}}$ 
}&
\parbox{0.88\textwidth}
{
    \vspace{0.2cm}
    Uncertainty in NC/CCQE cross section ratio: 36\% \cite{NIWGErr}.
    Applied to both corresponding SuperK and ND280 MC templates.\\   
}\\
\parbox{0.12\textwidth}
{ 
    \vspace{0.2cm}
    $f_{FSI;r}$ 
}&
\parbox{0.88\textwidth}
{
    \vspace{0.2cm}
    Uncertainty in effect of FSI (final state interaction / intranuclear hadron transport)
    to exclusive CCQE-like (1 $\mu$-like ring) SuperK measurement (e.g.\ effect of pion absorption in oxygen).
    Applied to SuperK MC templates.
    The error has dependence on reconstructed energy \cite{NIWGErr}.\\
}\\
\hline
\end{tabular}\\[2pt]
\caption{
  Inputs for $\nu_{\mu}$-disappearance systematic error evaluation (this is an updated version of Tab.\ 6 shown in Ref.\ \cite{OurRun1DisappAnalysis}).
}
\label{tab:input_uncertainties}
\end{table*}
